\chapter{Canalith Repositioning (Epley Maneuver)}

\section*{Overview}
Canalith repositioning, commonly known as the Epley maneuver, is a sequence of head and body movements that moves displaced calcium crystals out of the semicircular canals of the inner ear. It is the standard bedside treatment for benign paroxysmal positional vertigo (BPPV).
\section*{Alternative Names}
Epley maneuver; canalith repositioning procedure; particle repositioning maneuver; Dix-Hallpike-guided repositioning
\section*{Principle}
In BPPV, calcium carbonate debris (canaliths) from the utricle falls into a semicircular canal, usually the posterior canal, causing vertigo with head movement. Sequential head turns and positional changes use gravity to guide the debris back into the utricle, where it no longer stimulates the canal during movement.
\section*{History}
The Epley maneuver was described by John Epley in 1980. The associated Dix-Hallpike test, used to diagnose and position the patient, was introduced in 1952.
\section*{Why It Is Done}
Ordered for classic BPPV presenting as brief episodes of rotational vertigo triggered by rolling over in bed, looking up, or bending forward, confirmed by the Dix-Hallpike test.
\section*{Is This a Primary or Secondary Test or Procedure}
Primary definitive treatment for posterior canal BPPV, and used for other canal variants with modified positioning sequences.
\section*{How It Is Performed}
The patient sits upright with the head turned 45 degrees toward the affected ear, then lies back with the head hanging below the bed. The head is turned 90 degrees to the opposite side, the body is rolled onto the same side, and the patient finally sits up. Each position is held about 30 seconds. The sequence may be repeated until the provoking nystagmus resolves.
\section*{Preparation}
No specific preparation is required. A diagnosis of BPPV should be confirmed by the Dix-Hallpike test, and the maneuver is often performed in the clinic with the patient able to recline.
\section*{Contraindications and Precautions}
Contraindications include cervical spine instability, cervical radiculopathy, severe kyphosis, and high-grade carotid stenosis. Precautions include a history of recent neck trauma or surgery and limited neck mobility.
\section*{Risks}
Transient worsening of vertigo and nausea during the maneuver, vomiting, and rarely conversion of posterior canal BPPV to a horizontal canal variant. Falls may occur if dizziness persists after standing.
\section*{Aftercare and Recovery}
The patient is advised to avoid rapid head movements and extreme head extension for 24-48 hours and may sleep propped up. Symptoms usually resolve immediately but may recur.
\section*{Outcome and Follow-up}
Resolution of positional vertigo and a negative repeat Dix-Hallpike test indicate success. Persistent vertigo suggests canal conversion, incomplete repositioning, or an alternative diagnosis.
\section*{Expected Results}
A negative Dix-Hallpike test with no nystagmus or vertigo after repositioning is considered a normal, successful result.
\section*{Follow-up}
Patients are reassessed in days to weeks; recurrent BPPV may be treated by repeating the maneuver, and patients may be taught to perform the maneuver at home. Imaging and vestibular testing are reserved for atypical or refractory cases.
\section*{When to Seek Medical Care}
Follow the procedure team's specific instructions. Seek urgent medical assessment for severe or worsening pain, uncontrolled bleeding, fever or signs of infection, trouble breathing, chest pain, fainting, new weakness or confusion, inability to urinate, or any rapidly worsening symptom. The appropriate warning signs depend on the procedure and the patient's medical condition.

\section*{References}
\begin{enumerate}
  \item MedlinePlus. Benign positional vertigo. U.S. National Library of Medicine; 2025.
  \item Current Medical Diagnosis and Treatment. McGraw-Hill; 2025.
\end{enumerate}

\clearpage
